% try: edit the Python on the page, or change the style to friendly: the highlighting is redone as you type
\documentclass[11pt]{article}
\usepackage[margin=1in]{geometry}
\usepackage{minted}
\usemintedstyle{tango}
\setminted{fontsize=\small, linenos, frame=lines, framesep=6pt, numbersep=8pt}

\begin{document}
\begin{center}{\LARGE\bfseries Code, highlighted}\\[3pt]{\large by Pygments, inside the build}\end{center}

\section*{Python}
\begin{minted}{python}
def knuth_plass(words, width):
    """Break a paragraph into lines, minimising the total badness."""
    best = {0: (0, None)}
    for j in range(1, len(words) + 1):
        best[j] = min(
            (best[i][0] + badness(words[i:j], width), i)
            for i in range(j) if fits(words[i:j], width)
        )
    return lines_from(best, len(words))
\end{minted}

\section*{Rust}
\begin{minted}{rust}
/// The pages after an edit: rebuilt until they agree with the last run.
fn rebuild(cuts: &[Cut], edit: usize, last: &Run) -> Vec<Page> {
    let mut i = cuts.partition_point(|c| c.pos <= edit) - 1;
    let mut pages = Vec::new();
    loop {
        let (state, page) = typeset(&cuts[i], &cuts[i + 1]);
        pages.push(page);
        i += 1;
        if state == last.state_at(i) { break; }
    }
    pages
}
\end{minted}

\section*{C}
\begin{minted}{c}
#include <stdio.h>

static unsigned long factorial(unsigned n) {
    return n <= 1 ? 1 : n * factorial(n - 1);
}

int main(void) {
    for (int n = 1; n <= 5; n++)
        printf("%d! = %lu\n", n, factorial(n));
    return 0;
}
\end{minted}
\end{document}
